\documentclass[11pt]{article}
% Кодировка и язык
\usepackage[utf8]{inputenc}           % Кодировка исходного файла (UTF-8)
\usepackage[T2A]{fontenc}             % Кодировка шрифтов для кириллицы
\usepackage[english,russian]{babel}   % Локализация: русские названия, переносы

% Поиск в PDF
\usepackage{cmap}

% Поля страницы
\usepackage{geometry}
\geometry{left=3.8cm, right=3.8cm, top=2.5cm, bottom=2.5cm, paperheight=29.7cm}

% Математика
\usepackage{amsmath, amssymb, amsthm} % Базовые пакеты для математики
\usepackage{mathtools}                % Расширение amsmath (доп. символы и окружения)

% Картинки
\usepackage{graphicx}                 % Вставка картинок (\includegraphics)
\graphicspath{{pictures/}}            % Путь к папке с картинками
\usepackage{float}                    % Позволяет использовать [H] для жёсткого позиционирования
\usepackage{wrapfig}                  % Обтекание текста вокруг картинок

% Таблицы
\usepackage{array}                    % Расширенные возможности для таблиц
\usepackage{booktabs}                 % Красивые таблицы (\toprule, \midrule, \bottomrule)
\usepackage{multirow}                 % Объединение строк в таблицах
\usepackage{tabularx}                 % Таблицы с автоматической шириной колонок
\usepackage{ragged2e}                 % Улучшенное выравнивание (используется в tabularx)

% Списки
\usepackage{enumitem}                 % Гибкая настройка списков (itemize, enumerate)

% Колонтитулы
\usepackage{fancyhdr}

% Гиперссылки
\usepackage{xcolor}                   % Работа с цветами
\usepackage{hyperref}                 % Гиперссылки, оглавление, перекрёстные ссылки

% Умные ссылки
\usepackage{cleveref}

% Типографика
\usepackage{microtype}                % Улучшает переносы и межсловные расстояния

% Подписи к картинкам и таблицам
\usepackage{caption}                  % Настройка подписей (шрифт, отступы)

% Рисование схем и графиков
\usepackage{tikz}                     % Рисование схем, графиков, диаграмм
\usepackage{pgfplots}                 % Построение графиков функций (на основе tikz)
\pgfplotsset{compat=newest}           % Версия совместимости pgfplots

% Алгоритмы и код
\usepackage{algorithm}                % Окружения для алгоритмов
\usepackage{algorithmic}              % Синтаксис для алгоритмов
\usepackage{listings}                 % Вставка кода с подсветкой синтаксиса

\usepackage{soul}                     % Выделение текста: подчёркивание, зачёркивание, подсветка
\usepackage{multicol}                 % Многоколоночная вёрстка текста
\usepackage{titlesec}                 % Настройка внешнего вида заголовков секций
\usepackage{parskip}                  % Отступы между абзацами вместо красной строки
\usepackage{setspace}                 % Настройка межстрочного интервала
\usepackage{tcolorbox}                % Цветные блоки для заметок, цитат и выделений

% Палитра Quiet Light из VS Code
\definecolor{quietLightForeground}{HTML}{333333}
\definecolor{quietLightEditorBg}{HTML}{F5F5F5}
\definecolor{quietLightPanelBg}{HTML}{F2F2F2}
\definecolor{quietLightBorder}{HTML}{C9D0D9}
\definecolor{quietLightComment}{HTML}{AAAAAA}
\definecolor{quietLightKeyword}{HTML}{4B69C6}
\definecolor{quietLightType}{HTML}{7A3E9D}
\definecolor{quietLightConstant}{HTML}{9C5D27}
\definecolor{quietLightString}{HTML}{448C27}
\definecolor{quietLightFunction}{HTML}{AA3731}
\definecolor{quietLightPunctuation}{HTML}{777777}

% Совместимость со старыми именами цветов
\colorlet{pblue}{quietLightKeyword}
\colorlet{pgreen}{quietLightString}
\colorlet{pred}{quietLightFunction}
\colorlet{pgrey}{quietLightPunctuation}

\onehalfspacing

\titleformat{\section}
  {\Large\bfseries}
  {\thesection}
  {0.7em}
  {}

\titleformat{\subsection}
  {\large\bfseries}
  {\thesubsection}
  {0.7em}
  {}

\lstset{
    basicstyle=\ttfamily\small\color{quietLightForeground},
    backgroundcolor=\color{quietLightEditorBg},
    frame=single,
    rulecolor=\color{quietLightBorder},
    breaklines=true,
    columns=fullflexible,
    keepspaces=true,
    showstringspaces=false,
    keywordstyle=\color{quietLightKeyword}\bfseries,
    morekeywords=[2]{define,elif,else,endif,error,if,ifdef,ifndef,include,line,pragma,undef,warning},
    keywordstyle=[2]\color{quietLightKeyword},
    deletekeywords={bool,char,double,float,int,long,short,signed,unsigned,void},
    morekeywords=[3]{bool,char,double,float,int,long,short,signed,unsigned,void,size_t,string},
    keywordstyle=[3]\color{quietLightType},
    identifierstyle=\color{quietLightType},
    commentstyle=\color{quietLightComment}\itshape,
    stringstyle=\color{quietLightString},
    numberstyle=\color{quietLightConstant},
}

\lstdefinestyle{nolint}{
    basicstyle=\ttfamily\small\color{quietLightForeground},
    backgroundcolor=\color{quietLightPanelBg},
    columns=fullflexible,
    keepspaces=true,
    breaklines=false,
    keywordstyle=\color{quietLightKeyword},
    commentstyle=\color{quietLightKeyword},
    stringstyle=\color{quietLightKeyword},
    numberstyle=\color{quietLightKeyword},
}

\newcommand{\nolint}{\lstinline[style=nolint]}

\newtcolorbox{habrquote}{
    colback=quietLightPanelBg,
    colframe=quietLightBorder,
    boxrule=0.4pt,
    arc=2pt,
    left=8pt,
    right=8pt,
    top=6pt,
    bottom=6pt
}

% Настройка нумерации секций и подсекций
\renewcommand{\thesection}{\arabic{section}}
\renewcommand{\thesubsection}{\thesection.\arabic{subsection}}


% Настройка списков: убираем лишние отступы слева
\setlist[itemize]{leftmargin=*}
\setlist[enumerate]{leftmargin=*}


% Определение окружений для теорем, лемм, определений и т.д.
% Нумерация привязана к номеру секции
\newtheorem{theorem}{Теорема}[section]
\newtheorem{definition}{Определение}[section]
\newtheorem{remark}{Замечание}[section]
\newtheorem{lemma}{Лемма}[section]
\newtheorem{example}{Пример}[section]
\newtheorem{consequence}{Следствие}[section]
\newtheorem{algorithmm}{Алгоритм}
\newtheorem{question}{Вопрос}[section]
\newtheorem{answer}{Ответ}[section]
\newtheorem{statement}{Утверждение}[section]


\DeclareMathOperator{\sign}{sign}
\DeclareMathOperator{\Imm}{Im}
\DeclareMathOperator{\rg}{rg}

% Настройка колонтитулов
\pagestyle{fancy}
\fancyhf{} % Очистить все поля
\fancyhead[L]{Плюсистов поздравляем, остальным соболезнуем}
\fancyhead[R]{\href{https://t.me/ZNickA8871}{Автор}}
\fancyfoot[C]{\thepage}               % Справа внизу

\setlength{\headheight}{13.6pt}

\hypersetup{
    colorlinks=true,
    linkcolor=blue,
    citecolor=blue,
    filecolor=blue,
    urlcolor=blue,
}

\begin{document}

{\Huge\bfseries Лекция 1. Введение\par}

\vspace{2em}

\section{О чём предмет?}

Возьмём простой код:
\begin{lstlisting}[language=C++]
#include <iostream>

int main() {
  int x;
  std::cin >> x;
  std::cout << x + 5;
  return 0;
}
\end{lstlisting}
Скомпилируем его командой \nolint|g++ main.cpp|.

Получаем честный исполняемый файл в отличие от Java. Можем написать
\nolint|cat a.out| и получить белиберду (не совсем), а можем
\nolint|hexdump -c a.out| и получить бинарный формат (флагом включается отображение соответствующего байтам текста справа).

Вот мы вводим \nolint|5|, выводим \nolint|10|. Наша задача --- понять,
как и откуда это получилось. Вариант ``взяли из \nolint|iostream|'' нас
не устраивает: это мы должны знать с первого курса, сейчас нас интересует то, как наш ввод стал значениями переменных и как вычисленное значение попало в консоль.

\section{Этапы сборки}
\subsection{Препроцессинг}

Компиляция происходит в несколько стадий. Первая стадия --- препроцессинг.
Пишем \nolint|g++ -E [-v] main.cpp > main.e| и получаем результат стадии
препроцессинга. Флаг \nolint|-v| добавляет подробный вывод: по нему можно
посмотреть, откуда берутся библиотеки.

\begin{lstlisting}
#include "..." search starts here:
#include <...> search starts here:
 /usr/local/include
 .
 /usr/include/c++/14
 /usr/include/x86_64-linux-gnu/c++/14
 /usr/include/c++/14/backward
 /usr/lib/gcc/x86_64-linux-gnu/14/include
 /usr/include/x86_64-linux-gnu
 /usr/include
\end{lstlisting}

Идём в \nolint|iostream| и видим:

\begin{lstlisting}[language=C++]
extern istream cin;        ///< Linked to standard input
extern ostream cout;       ///< Linked to standard output
extern ostream cerr;       ///< Linked to standard error (unbuffered)
extern ostream clog;       ///< Linked to standard error (buffered)
\end{lstlisting}

Где-то на 60-х строчках лежат объявления переменных. Слово \nolint|extern| говорит нам о том, что эти переменные обязательно будут, но где-то там, далеко-далеко, а сейчас у нас есть только имя и тип (подробнее \href{https://en.cppreference.com/cpp/keyword/extern}{здесь}). Потом грепаем по
\nolint|main.e| и видим, что определения не появились.

Когда мы пишем \nolint|include|, мы подключаем только заголовочные файлы.
При сборке потом происходит линковка, в ходе которой подключаются уже собранные
реализации. На чистой системе исходного кода этих функций нет.

\subsection{Компиляция}

Теперь посмотрим следующую стадию --- компиляцию:

\begin{lstlisting}
g++ -S main.cpp
\end{lstlisting}

Команда \nolint|wc -l <file>| позволяет посчитать количество строк в файле.
Заметим, что здесь уже сильно меньше строк: порядка 120, тогда как в
\nolint|main.e| больше 30 тысяч.

Среди других инструкций найдём \nolint|call|. Это вызовы наших функций;
в таком виде они записываются в ассемблере. Здесь мы видим странные имена:
они изменены из-за того, что в C++ есть перегрузка функций, классы, операторы
и неймспейсы. Чтобы избежать дублирования имён, компилятор изменяет имена
таким образом. Это называется \textit{name mangling} или \textit{name decoration} (подробнее \href{https://en.wikipedia.org/wiki/Name_mangling}{здесь}).

Команда \nolint|c++filt <mangled name>| возвращает исходное имя объекта.

\subsection{Ассемблирование}

После этой стадии идёт ассемблирование, которое переводит файл из
\nolint|.s| в \nolint|.o|:

\begin{lstlisting}
g++ -c main.cpp
\end{lstlisting}

В нём инструкции просто отображаются в бинарный формат. Мы можем легко
отобразить объектный файл обратно в ассемблерный:

\begin{lstlisting}
objdump -d main.o
\end{lstlisting}

Но \nolint|main.o| мы всё ещё не можем исполнить. Почему? Потому что в нём
нет реализаций внешних функций.

\subsection{Линковка}

Соберём \nolint|main.cpp| полностью командой \nolint|g++ main.cpp| и
получим файл \nolint|a.out|. Напишем \nolint|objdump -d a.out|, найдём
\nolint|main| и наконец-то увидим рядом с \nolint|call| адрес.

Попросим подробный вывод (ключ -v) и увидим команду компиляции:

\begin{lstlisting}
/usr/libexec/gcc/x86_64-linux-gnu/14/cc1plus -quiet -v -imultiarch x86_64-linux-gnu -D_GNU_SOURCE main.cpp -quiet -dumpdir a- -dumpbase main.cpp -dumpbase-ext .cpp -mtune=generic -march=x86-64 -version -foffload-options=-l_GCC_m -fasynchronous-unwind-tables -fstack-protector-strong -Wformat -Wformat-security -fstack-clash-protection -fcf-protection -o /tmp/cch9iE9S.s
\end{lstlisting}

Затем команду ассемблирования:

\begin{lstlisting}
as -v --64 -o /tmp/ccpWsqla.o /tmp/cch9iE9S.s
\end{lstlisting}

И, наконец, команду линковки:

\begin{lstlisting}
/usr/libexec/gcc/x86_64-linux-gnu/14/collect2 -plugin /usr/libexec/gcc/x86_64-linux-gnu/14/liblto_plugin.so -plugin-opt=/usr/libexec/gcc/x86_64-linux-gnu/14/lto-wrapper -plugin-opt=-fresolution=/tmp/ccvUpgmp.res -plugin-opt=-pass-through=-lgcc_s -plugin-opt=-pass-through=-lgcc -plugin-opt=-pass-through=-lc -plugin-opt=-pass-through=-lgcc_s -plugin-opt=-pass-through=-lgcc --build-id --eh-frame-hdr -m elf_x86_64 --hash-style=gnu --as-needed -dynamic-linker /lib64/ld-linux-x86-64.so.2 -pie -z now -z relro /usr/lib/gcc/x86_64-linux-gnu/14/../../../x86_64-linux-gnu/Scrt1.o /usr/lib/gcc/x86_64-linux-gnu/14/../../../x86_64-linux-gnu/crti.o /usr/lib/gcc/x86_64-linux-gnu/14/crtbeginS.o -L/usr/lib/gcc/x86_64-linux-gnu/14 -L/usr/lib/gcc/x86_64-linux-gnu/14/../../../x86_64-linux-gnu -L/usr/lib/gcc/x86_64-linux-gnu/14/../../../../lib -L/lib/x86_64-linux-gnu -L/lib/../lib -L/usr/lib/x86_64-linux-gnu -L/usr/lib/../lib -L/usr/lib/gcc/x86_64-linux-gnu/14/../../.. /tmp/ccpWsqla.o -lstdc++ -lm -lgcc_s -lgcc -lc -lgcc_s -lgcc /usr/lib/gcc/x86_64-linux-gnu/14/crtendS.o /usr/lib/gcc/x86_64-linux-gnu/14/../../../x86_64-linux-gnu/crtn.o
\end{lstlisting}

\section{Исполняемые файлы и библиотеки}

\subsection{Динамические библиотеки}

Есть команда \nolint|ldd|, которая позволяет посмотреть динамические
библиотеки, необходимые для запуска бинаря:

\begin{lstlisting}
linux-vdso.so.1
libstdc++.so.6 => /lib/x86_64-linux-gnu/libstdc++.so.6
libc.so.6 => /lib/x86_64-linux-gnu/libc.so.6
libm.so.6 => /lib/x86_64-linux-gnu/libm.so.6
/lib64/ld-linux-x86-64.so.2
libgcc_s.so.1 => /lib/x86_64-linux-gnu/libgcc_s.so.1
\end{lstlisting}

Воспользуемся командой \nolint|file|:

\begin{lstlisting}
/lib/x86_64-linux-gnu/libstdc++.so.6: symbolic link to libstdc++.so.6.0.34
/lib/x86_64-linux-gnu/libstdc++.so.6.0.34: ELF 64-bit LSB shared object, x86-64, version 1 (GNU/Linux), dynamically linked, BuildID[sha1]=8d4f2235ec34ae33c412aa436c18ef4618f2efa6, stripped
a.out: ELF 64-bit LSB pie executable, x86-64, version 1 (SYSV), dynamically linked, interpreter /lib64/ld-linux-x86-64.so.2, BuildID[sha1]=ba1ae9a1317b2e65de6eab8d0c4c45f0c952dd77, for GNU/Linux 3.2.0, not stripped
\end{lstlisting}

Посмотрим, что они все формата ELF (\textit{Executable and Linkable Format}).
Это формат бинарных файлов во всех Unix-системах.

\subsection{ELF}

Есть тулза \nolint|readelf|, она позволяет посмотреть информацию о таких
файлах. Например:

\begin{lstlisting}
readelf -h a.out
\end{lstlisting}
Заголовок исполняемого файла:

\begin{lstlisting}
ELF Header:
  Magic:   7f 45 4c 46 02 01 01 00 00 00 00 00 00 00 00 00
  Class:                             ELF64
  Data:                              2's complement, little endian
  Version:                           1 (current)
  OS/ABI:                            UNIX - System V
  ABI Version:                       0
  Type:                              DYN (Position-Independent Executable file)
  Machine:                           Advanced Micro Devices X86-64
  Version:                           0x1
  Entry point address:               0x10a0
  Start of program headers:          64 (bytes into file)
  Start of section headers:          14472 (bytes into file)
  Flags:                             0x0
  Size of this header:               64 (bytes)
  Size of program headers:           56 (bytes)
  Number of program headers:         14
  Size of section headers:           64 (bytes)
  Number of section headers:         32
  Section header string table index: 31
\end{lstlisting}
\newpage

Заголовок динамической библиотеки:

\begin{lstlisting}
ELF Header:
  Magic:   7f 45 4c 46 02 01 01 03 00 00 00 00 00 00 00 00
  Class:                             ELF64
  Data:                              2's complement, little endian
  Version:                           1 (current)
  OS/ABI:                            UNIX - GNU
  ABI Version:                       0
  Type:                              DYN (Shared object file)
  Machine:                           Advanced Micro Devices X86-64
  Version:                           0x1
  Entry point address:               0x0
  Start of program headers:          64 (bytes into file)
  Start of section headers:          2651456 (bytes into file)
  Flags:                             0x0
  Size of this header:               64 (bytes)
  Size of program headers:           56 (bytes)
  Number of program headers:         13
  Size of section headers:           64 (bytes)
  Number of section headers:         35
  Section header string table index: 34
\end{lstlisting}

Те объекты, которыми оперирует линковщик, называются символами. Смотрим символы
файла командой \nolint|readelf -s|. Например, main.o:

\begin{lstlisting}
Symbol table '.symtab' contains 13 entries:
   Num:    Value          Size Type    Bind   Vis      Ndx Name
     0: 0000000000000000     0 NOTYPE  LOCAL  DEFAULT  UND
     1: 0000000000000000     0 FILE    LOCAL  DEFAULT  ABS main.cpp
     2: 0000000000000000     0 SECTION LOCAL  DEFAULT    1 .text
     3: 0000000000000000     1 OBJECT  LOCAL  DEFAULT    5 _ZNSt8__detail30[...]
     4: 0000000000000001     1 OBJECT  LOCAL  DEFAULT    5 _ZNSt8__detail30[...]
     5: 0000000000000002     1 OBJECT  LOCAL  DEFAULT    5 _ZNSt8__detail30[...]
     6: 0000000000000000     0 NOTYPE  GLOBAL DEFAULT  UND _ZSt21ios_base_l[...]
     7: 0000000000000000    99 FUNC    GLOBAL DEFAULT    1 main
     8: 0000000000000000     0 NOTYPE  GLOBAL DEFAULT  UND _ZSt3cin
     9: 0000000000000000     0 NOTYPE  GLOBAL DEFAULT  UND _ZNSirsERi
    10: 0000000000000000     0 NOTYPE  GLOBAL DEFAULT  UND _ZSt4cout
    11: 0000000000000000     0 NOTYPE  GLOBAL DEFAULT  UND _ZNSolsEi
    12: 0000000000000000     0 NOTYPE  GLOBAL DEFAULT  UND __stack_chk_fail
\end{lstlisting}

\newpage

Посмотрим, где лежит \nolint|cin|:

\begin{lstlisting}
readelf -s /lib/x86_64-linux-gnu/libstdc++.so.6 | grep St3cin
\end{lstlisting}

\begin{lstlisting}
2306: 00000000002896a0   280 OBJECT  GLOBAL DEFAULT   31 _ZSt3cin@@GLIBCXX_3.4
\end{lstlisting}

А вот так можно посмотреть количество символов в \nolint|libstdc++|:

\begin{lstlisting}
readelf -s /lib/x86_64-linux-gnu/libstdc++.so.6 | wc -l
\end{lstlisting}

Мы можем переименовать \nolint|libstdc++| (главное не забыть вернуть всё как
было) и увидеть, что наш ``самодостаточный'' бинарник не запустится. Всё потому,
что он зависит от динамических библиотек.

С помощью \nolint|readelf -a a.out| можно посмотреть всю
информацию о файле. А переменной окружения \nolint|LD_PRELOAD| можно указать
другую версию библиотеки:

\begin{lstlisting}
LD_PRELOAD=/lib/x86_64-linux-gnu/libstdc++.so.6.0.33 ./a.out
\end{lstlisting}

Если включить статическую линковку ( это делается флагом \nolint|-static|), заголовок будет выглядеть иначе:

\begin{lstlisting}
ELF Header:
  Magic:   7f 45 4c 46 02 01 01 03 00 00 00 00 00 00 00 00
  Class:                             ELF64
  Data:                              2's complement, little endian
  Version:                           1 (current)
  OS/ABI:                            UNIX - GNU
  ABI Version:                       0
  Type:                              EXEC (Executable file)
  Machine:                           Advanced Micro Devices X86-64
  Version:                           0x1
  Entry point address:               0x405e20
  Start of program headers:          64 (bytes into file)
  Start of section headers:          2415752 (bytes into file)
  Flags:                             0x0
  Size of this header:               64 (bytes)
  Size of program headers:           56 (bytes)
  Number of program headers:         11
  Size of section headers:           64 (bytes)
  Number of section headers:         29
  Section header string table index: 28
\end{lstlisting}

\end{document}
